\documentclass[11pt]{article}

\usepackage[margin=1in]{geometry}
\usepackage{graphicx}
\usepackage[colorlinks=true,allcolors=blue]{hyperref}
\usepackage{authblk}
\usepackage[style=numeric, sorting=none, backend=biber]{biblatex}
\addbibresource{references.bib}



% determines the spacing and general formatting of each section in the document
\makeatletter
\renewcommand{\@maketitle}{%
    \newpage
    \null
    \vskip 2em%
    \let\footnote\thanks
    \begin{center}%
        {\LARGE \@title \par}%
    \end{center}%
        \vskip 1em%
        {\large \noindent\begin{minipage}{\textwidth}\raggedright \@author\end{minipage}\par}%
        \vskip 1em%
    \begin{center}%
        {\@date}%
    \end{center}%
    \par
    \vskip 1.5em%
}

% Compact bibliography format: [N] Author et al., Year, ShortJournal.
\DeclareNameFormat{familyonly}{%
  \ifnumequal{\value{listcount}}{1}{%
    \namepartfamily%
    \ifnumgreater{\value{listtotal}}{1}{ et al.}{}%
  }{}%
}
\DeclareBibliographyDriver{article}{%
  \printnames[familyonly]{author}%
  , \printfield{year}%
  , \printfield{shortjournal}%
  }
\DeclareBibliographyDriver{inproceedings}{%
  \printnames[familyonly]{author}%
  , \printfield{year}%
  , \printfield{shortjournal}%
}
\AtBeginBibliography{\setlength{\parskip}{0pt}}
\setlength{\bibitemsep}{0pt}
\setlength{\bibparsep}{0pt}

\makeatother

\title{Constraining the Viscosity of Europa's Ice Shell from Eutectic Interfaces in Geodynamic Models}
\renewcommand\Authands{, }
\author[1,2]{William Byrne}
\author[1]{Ana-Catalina Plesa}
\author[1]{Hauke Hussmann}
\author[3]{Dustin Schroeder}
\author[4]{Gregor Steinbr\"ugge}
\author[5]{Natalie Wolfenbarger}

\affil[1]{Institute of Space Research, German Aerospace Center (DLR), Berlin, Germany\\
\href{mailto:William.Byrne@dlr.de}{William.Byrne@dlr.de}; \href{mailto:Ana.Plesa@dlr.de}{Ana.Plesa@dlr.de}; \href{mailto:Hauke.Hussmann@dlr.de}{Hauke.Hussmann@dlr.de}; \href{mailto:gael.cascioli@dlr.de}{gael.cascioli@dlr.de}}
\affil[2]{Freie Universit\"at Berlin, Berlin, Germany\\
\href{mailto:byrnew02@zedat.fu-berlin.de}{byrnew02@zedat.fu-berlin.de}; \href{mailto:frank.postberg@fu-berlin.de}{frank.postberg@fu-berlin.de}}
\affil[3]{Department of Geophysics, Stanford University, Stanford, CA, USA\\
\href{mailto:Dustin.M.Schroeder@stanford.edu}{Dustin.M.Schroeder@stanford.edu}}
\affil[4]{Jet Propulsion Laboratory, California Institute of Technology, Pasadena, CA, USA\\
\href{mailto:gregor.b.steinbruegge@nasa.jpl.gov}{gregor.b.steinbruegge@nasa.jpl.gov}}
\affil[5]{Los Alamos National Laboratory, Los Alamos, NM, USA\\
\href{mailto:wolfenbarger@lanl.gov}{wolfenbarger@lanl.gov}}
\date{21.01.2026:\href{https://meetingorganizer.copernicus.org/EGU25/EGU25-652.html}{View abstract on the EGU2025 website}}


\begin{document}

\maketitle

\raggedright
\setlength{\parindent}{0pt}
\setlength{\parskip}{\baselineskip}

\section*{Abstract}

Jupiter’s moon Europa is one of the prime targets for planetary exploration due to its high astrobiological potential. Slightly smaller than Earth’s moon, Europa harbors a liquid water ocean beneath an ice shell. The thickness of Europa’s ice shell is poorly constrained and values of less than 1 km  to up to 90 km have been suggested in previous studies \parencite{billingsGreatThicknessDebate2005, vilellaTidallyHeatedConvection2020}. Ice-penetrating radars on NASA’s Europa Clipper (REASON, \parencite{blankenship2009radar}) and ESA’s JUICE (RIME, \parencite{bruzzoneRIMERadarIcy2013a}) missions aim to determine the thickness of Europa's ice shell. Recent studies have suggested that constraints on the thickness of Europa’s ice shell can be obtained through the detection of eutectic interfaces, defined as the depth where brine becomes thermodynamically stable in the ice shell \parencite{schroederConstrainingThicknessConductive2024}. In fact, previous studies have shown that the detection of eutectic horizons within an ice shell is likely easier than detecting the ice-ocean interface, given their shallower depths and therefore lower total signal attenuation \parencite{kalousovaRadarAttenuationEuropas2017,soucekRadarAttenuationEnceladus2023,byrneRadarAttenuationSubsurface2024}. The depth of the eutectic interfaces depends on the thermal state of the ice shell, which is closely linked to the ice shell viscosity and large-scale dynamics \parencite{kalousovaRadarAttenuationEuropas2017}. As suggested by previous authors \parencite{kalousovaRadarAttenuationEuropas2017, schroederConstrainingThicknessConductive2024}, detection of eutectic interfaces therefore represents a promising strategy to constrain the thermophysical properties of the ice shell through characterization of its convective pattern.

In this study we use the geodynamic code GAIA \parencite{huttigImprovedFormulationIncompressible2013} to investigate the ice shell dynamics on Europa. We vary the ice shell thickness and ice shell viscosity that largely affect the convection pattern and in particular the number of hot upwellings and cold downwellings that can develop. In our models, the viscosity is temperature dependent and follows an Arrhenius law. We choose a reference value for the viscosity at the ice-ocean interface and vary this over several orders of magnitude between the different models. Once a simulation has reached a statistical (quasi-)steady state, we determine the eutectic pattern by identifying the depths of the eutectic temperature. We treat this sequence of eutectic depths as a signal and identify the peaks of each local maxima (or peak) in the signal. The number of local maxima in the simulation is used to estimate the global number of convection cells in the ice shell. 

Our preliminary results show a close relation between the number of plumes that develop in the ice shell of Europa and the viscosity at the ice-ocean interface. By increasing the number and complexity of our simulations, we aim to derive so-called scaling laws that will relate the convection structure with the viscosity and thickness of Europa’s ice shell. This will provide a framework that will help to interpret the detection of eutectic interfaces in future radar measurements in the context of large-scale dynamics of the deep ice shell. 

\printbibliography[title={References}]

\end{document}
